% Options for packages loaded elsewhere
\PassOptionsToPackage{unicode}{hyperref}
\PassOptionsToPackage{hyphens}{url}
\documentclass[
  ignorenonframetext,
  aspectratio=169,
]{beamer}
\newif\ifbibliography
\usepackage{pgfpages}
\setbeamertemplate{caption}[numbered]
\setbeamertemplate{caption label separator}{: }
\setbeamercolor{caption name}{fg=normal text.fg}
\beamertemplatenavigationsymbolsempty
% remove section numbering
\setbeamertemplate{part page}{
  \centering
  \begin{beamercolorbox}[sep=16pt,center]{part title}
    \usebeamerfont{part title}\insertpart\par
  \end{beamercolorbox}
}
\setbeamertemplate{section page}{
  \centering
  \begin{beamercolorbox}[sep=12pt,center]{section title}
    \usebeamerfont{section title}\insertsection\par
  \end{beamercolorbox}
}
\setbeamertemplate{subsection page}{
  \centering
  \begin{beamercolorbox}[sep=8pt,center]{subsection title}
    \usebeamerfont{subsection title}\insertsubsection\par
  \end{beamercolorbox}
}
% Prevent slide breaks in the middle of a paragraph
\widowpenalties 1 10000
\raggedbottom
\AtBeginPart{
  \frame{\partpage}
}
\AtBeginSection{
  \ifbibliography
  \else
    \frame{\sectionpage}
  \fi
}
\AtBeginSubsection{
  \frame{\subsectionpage}
}
\usepackage{iftex}
\ifPDFTeX
  \usepackage[T1]{fontenc}
  \usepackage[utf8]{inputenc}
  \usepackage{textcomp} % provide euro and other symbols
\else % if luatex or xetex
  \usepackage{unicode-math} % this also loads fontspec
  \defaultfontfeatures{Scale=MatchLowercase}
  \defaultfontfeatures[\rmfamily]{Ligatures=TeX,Scale=1}
\fi
\usepackage{lmodern}
\ifPDFTeX\else
  % xetex/luatex font selection
\fi
% Use upquote if available, for straight quotes in verbatim environments
\IfFileExists{upquote.sty}{\usepackage{upquote}}{}
\IfFileExists{microtype.sty}{% use microtype if available
  \usepackage[]{microtype}
  \UseMicrotypeSet[protrusion]{basicmath} % disable protrusion for tt fonts
}{}
\makeatletter
\@ifundefined{KOMAClassName}{% if non-KOMA class
  \IfFileExists{parskip.sty}{%
    \usepackage{parskip}
  }{% else
    \setlength{\parindent}{0pt}
    \setlength{\parskip}{6pt plus 2pt minus 1pt}}
}{% if KOMA class
  \KOMAoptions{parskip=half}}
\makeatother
\usepackage{longtable,booktabs,array}
\usepackage{caption}
\captionsetup[table]{skip=6pt}
\newcounter{none} % for unnumbered tables
\usepackage{calc} % for calculating minipage widths
\usepackage{caption}
% Make caption package work with longtable
\makeatletter
\def\fnum@table{\tablename~\thetable}
\makeatother
\setlength{\emergencystretch}{3em} % prevent overfull lines
\providecommand{\tightlist}{%
  \setlength{\itemsep}{0pt}\setlength{\parskip}{0pt}}
% Auto-generated by infrastructure/rendering/_pdf_latex_helpers.py::extract_math_font_preamble.
% Minimal header passed to Pandoc via -H so Beamer slides pick up the same math font
% as the combined-PDF path. Adjust the manuscript's preamble.md to change the font globally.
\usepackage{unicode-math}
\setmathfont{latinmodern-math.otf}

% Manuscript macro/package fallbacks (newcommand -> providecommand).
% Core mathematics
\usepackage{amsmath}
\usepackage{amssymb}
% Document layout
% Tables
\usepackage{booktabs}
% Caption styling (booktabs already loaded above). Small caption font with a
% bold label ("Figure 1", "Table 2") gives the math/figure-heavy layout a
% consistent, journal-like caption rhythm without fighting pandoc-crossref —
% pandoc emits standard \caption{}, and caption/captionsetup only restyle it.
% ── Theorem environments ─────────────────────────────────────────────────
% amsthm is loaded above. The FedGVI<->Friston bridge is stated as a small
% number of theorems/definitions/lemmas/propositions/corollaries; sharing one
% counter (definition/lemma/proposition/corollary all step the `theorem`
% counter) gives a single monotone numbering 1, 2, 3, ... across all kinds, so
% authors can reference them as Theorem 1, Definition 2, Corollary 3 without a
% per-kind counter drifting out of order. Plain style for theorem-like results,
% definition style (upright body) for definitions.
% (algorithm2e intentionally not loaded: this manuscript states results as
% numbered amsthm Theorems/Definitions, not algorithm2e pseudocode.)
% Code listings
\usepackage{listings}
% Typography and formatting
% (siunitx intentionally not loaded: no \SI/\num units appear in this manuscript.)
% Cross-references and citations
% (cleveref and titlesec intentionally not loaded: unavailable in this TeX tree
% and not required. Cross-references resolve through pandoc-crossref's own
% \ref-style names — no \cref appears — and section heading styling falls back
% to the pandoc/LaTeX defaults.)
% ── TeX Live Basic-compatible mono font for code listings ────────────
% Latin Modern Mono is bundled with TeX Live Basic and keeps the manuscript
% renderable on minimal TeX installations. Mathematical Unicode coverage is
% provided separately by Latin Modern Math below, so the code font does not
% need a private JuliaMono installation.
%
% Use the TeX font filename rather than the system-family name: the minimal
% TeX Live fontconfig database may not register the OpenType family even though
% kpsewhich can resolve the bundled files.
% Math font for unicode-math: Latin Modern Math (TeX Live) has full BMP
% coverage including U+2223 (\mid), U+226A/226B (\ll/\gg), and the Greek/
% blackboard letters used in equations. Without an explicit \setmathfont,
% unicode-math falls back to lmroman text font which lacks several glyphs
% and emits "Missing character" warnings on every \mid in math mode.

% Slide layout defaults for warning-clean scientific decks.
\usepackage{etoolbox}
\IfFileExists{xurl.sty}{\usepackage{xurl}}{}
\IfFileExists{seqsplit.sty}{\usepackage{seqsplit}}{\newcommand{\seqsplit}[1]{#1}}
\protected\def\breakseq#1{\seqsplit{#1}}
\protected\def\breaktt#1{\begingroup\ttfamily\seqsplit{#1}\endgroup}
\setlength{\emergencystretch}{6em}
\tolerance=5000
\hbadness=10000
\hfuzz=1pt
\setlength{\tabcolsep}{2pt}
\AtBeginEnvironment{longtable}{\tiny\renewcommand{\arraystretch}{0.86}\setlength{\tabcolsep}{1pt}}
\AtBeginEnvironment{tabular}{\tiny\renewcommand{\arraystretch}{0.86}\setlength{\tabcolsep}{1pt}}
\AtBeginEnvironment{equation}{\tiny}
\AtBeginEnvironment{equation*}{\tiny}
\AtBeginEnvironment{align}{\tiny}
\AtBeginEnvironment{align*}{\tiny}
\AtBeginEnvironment{itemize}{\footnotesize}
\AtBeginEnvironment{enumerate}{\footnotesize}
\AtBeginEnvironment{description}{\footnotesize}
\setbeamerfont{caption}{size=\tiny}
\setbeamerfont{caption name}{size=\tiny}
\setbeamerfont{normal text}{size=\small}
\setbeamerfont{frametitle}{size=\small}
\setbeamerfont{section title}{size=\footnotesize}
\setbeamerfont{subsection title}{size=\footnotesize}
\setbeamertemplate{section page}{%
  \centering
  \begin{beamercolorbox}[sep=12pt,center,wd=\paperwidth]{section title}
    \parbox{0.86\paperwidth}{\centering\usebeamerfont{section title}\insertsection\par}
  \end{beamercolorbox}
}
\setbeamertemplate{subsection page}{%
  \centering
  \begin{beamercolorbox}[sep=8pt,center,wd=\paperwidth]{subsection title}
    \parbox{0.86\paperwidth}{\centering\usebeamerfont{subsection title}\insertsubsection\par}
  \end{beamercolorbox}
}
\setlength{\abovecaptionskip}{2pt}
\setlength{\belowcaptionskip}{0pt}

% Natbib and cross-reference fallbacks — slides don't load natbib
% or cleveref, but combined-PDF manuscript prose may emit these
% commands. The fallback renders citations as a bracketed key list
% and cross-references as detokenized labels so slides don't fail on
% undefined control sequences or raw underscores. \providecommand is
% a no-op if packages load later.
\providecommand{\citep}[1]{[#1]}
\providecommand{\citet}[1]{#1}
\providecommand{\citealp}[1]{#1}
\providecommand{\citeauthor}[1]{#1}
\providecommand{\citeyear}[1]{#1}
\providecommand{\cref}[1]{\texttt{\detokenize{#1}}}
\providecommand{\Cref}[1]{\texttt{\detokenize{#1}}}
\providecommand{\crefrange}[2]{\texttt{\detokenize{#1}}--\texttt{\detokenize{#2}}}
\providecommand{\Crefrange}[2]{\texttt{\detokenize{#1}}--\texttt{\detokenize{#2}}}
\providecommand{\paragraph}[1]{\textbf{#1}\ }

% Opt-in accessible presentation profile.
\setbeamerfont{normal text}{size*={20pt}{24pt}}
\setbeamerfont{frametitle}{size*={28pt}{32pt}}
\setbeamerfont{section title}{size*={28pt}{32pt}}
\setbeamerfont{subsection title}{size*={28pt}{32pt}}
\setbeamerfont{caption}{size*={16pt}{19pt}}
\setbeamerfont{caption name}{size*={16pt}{19pt}}
\setbeamerfont{itemize/enumerate body}{size*={20pt}{24pt}}
\setbeamerfont{itemize/enumerate subbody}{size*={20pt}{24pt}}
\setbeamerfont{itemize/enumerate subsubbody}{size*={20pt}{24pt}}
\setbeamerfont{itemize item}{size*={20pt}{24pt}}
\setbeamerfont{itemize subitem}{size*={20pt}{24pt}}
\setbeamerfont{itemize subsubitem}{size*={20pt}{24pt}}
\setbeamerfont{description body}{size*={20pt}{24pt}}
\setbeamerfont{description item}{size*={20pt}{24pt}}
\setbeamerfont{quote}{size*={20pt}{24pt}}
\setbeamerfont{quotation}{size*={20pt}{24pt}}
\AtBeginEnvironment{longtable}{\fontsize{20pt}{24pt}\selectfont}
\AtBeginEnvironment{tabular}{\fontsize{20pt}{24pt}\selectfont}
\AtBeginEnvironment{equation}{\fontsize{20pt}{24pt}\selectfont}
\AtBeginEnvironment{equation*}{\fontsize{20pt}{24pt}\selectfont}
\AtBeginEnvironment{align}{\fontsize{20pt}{24pt}\selectfont}
\AtBeginEnvironment{align*}{\fontsize{20pt}{24pt}\selectfont}
\AtBeginEnvironment{itemize}{\fontsize{20pt}{24pt}\selectfont}
\AtBeginEnvironment{enumerate}{\fontsize{20pt}{24pt}\selectfont}
\AtBeginEnvironment{description}{\fontsize{20pt}{24pt}\selectfont}
\AtBeginDocument{\usebeamerfont{normal text}}
\newlength{\AFSlideTableWidth}
% The fixed footer reservation is calibrated with the semantic seven-line regular-body budget.
\setbeamertemplate{footline}{%
  \leavevmode\vbox{%
    \hbox{%
      \begin{beamercolorbox}[wd=\paperwidth,ht=16pt,dp=3pt,center]{author in head/foot}%
        \fontsize{16pt}{19pt}\selectfont Untagged PDF derivative \textbar\ \href{../web/index.html}{HTML reader}%
      \end{beamercolorbox}%
    }%
    \vskip6pt%
  }%
}
% Accessible footer reserved height: 25pt.

\newtheorem{proposition}{Proposition}
\newtheorem{hypothesis}{Hypothesis}
\newtheorem{remark}[theorem]{Remark}
\newtheorem{axiom}{Axiom}
\newtheorem{property}{Property}
\usepackage{bookmark}
\IfFileExists{xurl.sty}{\usepackage{xurl}}{} % add URL line breaks if available
\urlstyle{same}
% fallback for those not using the hyperref driver hyperxmp:
\makeatletter
\@ifundefined{xmpquote}{\newcommand{\xmpquote}[1]{#1}}{}
\makeatother
\hypersetup{
  hidelinks,
  pdfcreator={LaTeX via pandoc}}

\author{\texorpdfstring{}{}}
\date{}

\begin{document}

\begin{frame}{Related work: active inference, federated Bayes, and the
scoped bridge}
\protect\phantomsection\label{sec:related-work}
\end{frame}

\begin{frame}{Related work (part 2)}
This work sits at the boundary between two research communities with
different objects of inference. Each cited thread contributes a real
component; this paper extends the intersection without claiming that
either literature is exhausted. The positioning is therefore against the
reviewed sources and their assumptions,
\end{frame}

\begin{frame}{Related work (part 3)}
not against an absolute claim about everything the field has or has not
done.
\end{frame}

\begin{frame}{Pre-modern probability, inverse probability, and
collective judgment}
\protect\phantomsection\label{sec:related-historical}
These historical sources provide conceptual lineage, not evidence that
early authors anticipated KL minimization, active inference, or
federation. Pascal and Fermat, Huygens, Montmort, Bernoulli,
\end{frame}

\begin{frame}{Pre-modern probability\ldots{} (part 2)}
and de Moivre made expectation and uncertain evidence objects of
calculation (Pascal and Fermat 1654; Huygens 1657; Montmort 1708;
Bernoulli 1713; Moivre 1718).
\end{frame}

\begin{frame}{Pre-modern probability\ldots{} (part 3)}
Bayes and Price framed inference from observed events to an unknown
chance, and Laplace generalized inference about causes (Bayes 1763;
Laplace 1774). Their relevance here is vocabulary for generative-model
inversion. The executable recovery identity is instead a modern,
\end{frame}

\begin{frame}{Pre-modern probability\ldots{} (part 4)}
project-local result of sec.~5.8 and
sec.~6.
\end{frame}

\begin{frame}{Pre-modern probability\ldots{} (part 5)}
Daniel Bernoulli's expected-utility treatment and the aggregation work
of Borda and Condorcet connect uncertain judgment to action and
collective choice (Bernoulli 1738; Borda 1784; Condorcet 1785). They
motivate questions about competence, independence, weights, and stakes.
\end{frame}

\begin{frame}{Pre-modern probability\ldots{} (part 6)}
They do not formally support this paper's log-pool, generalized-Bayes,
or robustness claims; the full genealogy remains documented in the
source audit.
\end{frame}

\begin{frame}{Active inference: generative agents, EFE, and colonies}
\protect\phantomsection\label{sec:related-aif}
The free-energy principle (Friston 2010) and discrete state-space
process theory (Friston et al. 2017) connect generative models,
variational inference, and expected-free-energy action selection.
\end{frame}

\begin{frame}{Active inference (part 2)}
The discrete synthesis (Da Costa et al. 2020), pymdp (Heins et al.
2022), and RxInfer (Bagaev et al. 2023) provide executable machinery.
\end{frame}

\begin{frame}{Active inference (part 3)}
This project reimplements the categorical substrate and its
risk--ambiguity decomposition (eq.~22,
eq.~23, fig.~8), then evaluates
explicit robustness controls at belief fusion.
\end{frame}

\begin{frame}{Active inference (part 4)}
The fusion operator itself is not new as a mathematical object. Under
the explicit categorical posterior-log-potential and fixed-weight bridge
of sec.~5.8,
\end{frame}

\begin{frame}{Active inference (part 5)}
the message-combination term of Friston's belief-sharing equation is
represented by a logarithmic opinion pool (Genest and Zidek 1986; Genest
et al. 1986) and a product-of-experts consensus (Hinton 2002).
\end{frame}

\begin{frame}{Active inference (part 6)}
That limited representation is not a statement that the complete source
protocol is a log pool. Abbas' KL view of linear and log-linear pools
makes the same point in scoring-rule language: log-pooling can be
justified as a KL aggregation rule for expert distributions, not as a
contamination-robust estimator (Abbas 2009).
\end{frame}

\begin{frame}{Active inference (part 7)}
The Bayesian Committee Machine adds the distributed-learning analogue:
independent estimators trained on data subsets can be combined by a
product-style Bayesian rule, with assumptions about conditional
independence and prior accounting made explicit (Tresp 2000).
\end{frame}

\begin{frame}{Active inference (part 8)}
Fully Bayesian aggregation gives the social-choice counterpart:
geometric pooling of beliefs is singled out by dynamic Bayesian
rationality conditions, not by robustness to contaminated reports
(Dietrich 2021).
\end{frame}

\begin{frame}{Active inference (part 9)}
That classical literature is useful precisely because it names the
hidden commitments --- shared support, weight choice, prior accounting,
and external-Bayes coherence --- that are easy to overlook when the same
operation appears as a colony update.
\end{frame}

\begin{frame}{Active inference (part 10)}
A parallel thread studies how active-inference agents coordinate as
ensembles --- collective behavior and surprise minimization across a
colony (Heins et al. 2024), epistemic communities (Albarracin et al.
2022), and collective intelligence (Kaufmann et al. 2021).
\end{frame}

\begin{frame}{Active inference (part 11)}
The belief-sharing thread (Friston et al. 2024) is one member of this
family. Agents fuse beliefs, and our colony scenario is a reduced
categorical standard- Bayes-limit analogue of its worked example
(eq.~14, sec.~7.2).
\end{frame}

\begin{frame}{Active inference (part 12)}
Structure learning completes the conceptual frame. Active inference with
Bayesian optimal design selects and reduces models (Smith et al. 2022),
while post-hoc BMR (Friston and Penny 2011) motivates our configured
sign control (eq.~20,
sec.~7.4).
\end{frame}

\begin{frame}{Active inference (part 13)}
Among the active-inference sources reviewed here, belief fusion is not
systematically characterized under explicit contamination or
intentionally wrong broadcasts, nor connected to the robustness theory
used here. Fusion is treated as trusting in that scoped comparison.
\end{frame}

\begin{frame}{Active inference (part 14)}
Friston et al. (2024) present communicating-colony free-energy
convergence, Dirichlet language acquisition, and Bayesian model
reduction. We use reduced categorical standard-Bayes-limit analogues of
those mechanisms (sec.~7.1) before the robust
extension.
\end{frame}

\begin{frame}{Active inference (part 15)}
The third is a configured sign control on a fixed posterior. None is an
exact source-protocol or figure reproduction.
\end{frame}

\begin{frame}{Robust and federated Bayes outside active inference}
\protect\phantomsection\label{sec:related-fl}
Federated learning aggregates models trained on decentralized data
(McMahan et al. 2017), and the probabilistic-federation line recasts
that aggregation as variational inference ---
\end{frame}

\begin{frame}{Robust and federated Bayes\ldots{} (part 2)}
partitioned variational inference (Ashman et al. 2022) and its federated
predecessor (Bui et al. 2018). This is a different use of the word
\emph{federated} from Friston et al.'s federated inference:
\end{frame}

\begin{frame}{Robust and federated Bayes\ldots{} (part 3)}
Friston federates hidden-state beliefs among agents sharing a world
model, while machine-learning federated learning usually federates
parameter or predictive-model updates over decentralized datasets.
\end{frame}

\begin{frame}{Robust and federated Bayes\ldots{} (part 4)}
The bridge claimed here is therefore algebraic and variational --- a
shared normalized product/log-pool operator at the KL/NLL recovery
corner --- not a claim that either source paper already solved the
other's problem.
\end{frame}

\begin{frame}{Robust and federated Bayes\ldots{} (part 5)}
Robustness, meanwhile, has a mature Bayesian theory: general Bayesian
updating through a loss (Bissiri et al. 2016), Gibbs posterior inference
(Jiang and Tanner 2008), safe learning-rate selection under
misspecification (Grünwald 2012),
\end{frame}

\begin{frame}{Robust and federated Bayes\ldots{} (part 6)}
coarsened posteriors for robustness to exact-data conditioning (Miller
and Dunson 2018), divergence-criteria posterior updating (Jewson et al.
2018), Bayesian misspecification asymptotics (Kleijn and Vaart 2012).
\end{frame}

\begin{frame}{Robust and federated Bayes\ldots{} (part 7)}
The optimization-centric GVI view (Knoblauch et al. 2022) and recent
closed-form characterizations (Nguyen and Westerhout 2026) connect
objectives to approximating families.
\end{frame}

\begin{frame}{Robust and federated Bayes\ldots{} (part 8)}
Bounded-influence approaches include density-power and gamma-divergence
estimation (Basu et al. 1998; Fujisawa and Eguchi 2008; Ghosh and Basu
2015), robust-divergence variational inference (Futami et al. 2018), and
generalized cross-entropy (Zhang and Sabuncu 2018).
\end{frame}

\begin{frame}{Robust and federated Bayes\ldots{} (part 9)}
Huber and Ronchetti supply the robust-statistics vocabulary of influence
and breakdown (Huber and Ronchetti 2009), used here for boundedness and
empirical failure modes rather than theorem transfer.
\end{frame}

\begin{frame}{Robust and federated Bayes\ldots{} (part 10)}
FedGVI (Mildner, Hamelijnck, et al. 2025) supplies the per-agent
generalized-Bayes synthesis with client and server divergence choices.
\end{frame}

\begin{frame}{Robust and federated Bayes\ldots{} (part 11)}
This literature provides decentralized aggregation and theorem-backed
results in its stated settings. Among the sources reviewed here, it does
not evaluate active-inference POMDP belief consensus---the
generative-model-bearing, action-selecting setting where beliefs drive
behavior.
\end{frame}

\begin{frame}{Robust and federated Bayes\ldots{} (part 12)}
The 2025 preprint on convergence rates under prior misspecification
(Mildner, Giampouras, et al. 2025) sharpens the current GVI context:
bounded divergences can support concentration and rates under explicit
assumptions,
\end{frame}

\begin{frame}{Robust and federated Bayes\ldots{} (part 13)}
but the result does not transfer automatically to this repository's
finite categorical state space or to its server-side aggregation
heuristic.
\end{frame}

\begin{frame}{Robust and federated Bayes\ldots{} (part 14)}
The federated-learning robustness literature also supplies important
negative space for this paper. Byzantine-tolerant gradient aggregation
(Blanchard et al. 2017), geometric-median robust aggregation (Pillutla
et al. 2022),
\end{frame}

\begin{frame}{Robust and federated Bayes\ldots{} (part 15)}
and divergence-weighted gamma-mean aggregation (Li et al. 2022) attack
corrupted client updates directly. A recent Bayesian robust-aggregation
preprint likewise models unknown client honesty for federated model
updates (Karakulev et al. 2025).
\end{frame}

\begin{frame}{Robust and federated Bayes\ldots{} (part 16)}
Those are close comparators for adversarial federation, but the object
being aggregated differs: they aggregate model-update vectors or
posterior measures, while this paper attacks categorical belief fusion
and generalized-Bayes client updates.
\end{frame}

\begin{frame}{Robust and federated Bayes\ldots{} (part 17)}
Robust subset-posterior combination is another nearby Bayesian route
(Minsker et al. 2017), but it combines posterior measures across data
shards rather than active-inference belief broadcasts with a shared
latent state.
\end{frame}

\begin{frame}[fragile]{Robust and federated Bayes\ldots{} (part 18)}
We use those sources to position the problem, not to import their
guarantees into \texttt{robust\_aggregate}.
\end{frame}

\begin{frame}{The specific bridge added here}
\protect\phantomsection\label{sec:related-gap}
The scoped bridge evaluated by this manuscript is generalized-Bayes
belief fusion for active-inference ensembles. Concretely:
\end{frame}

\begin{frame}{The specific bridge added\ldots{} (part 2)}
\textbf{Per-agent bounded-loss updates.} FedGVI-faithful \(\beta\)/rcce
updates carry the cited bounded-influence result only under its source
assumptions (eq.~7, eq.~8).
\end{frame}

\begin{frame}{The specific bridge added\ldots{} (part 3)}
\textbf{Server-side divergence reweighting.} The complementary heuristic
has a stated recovery limit (eq.~10). Its finite
empirical behavior does not borrow theorems from robust federated
aggregation.
\end{frame}

\begin{frame}{The specific bridge added\ldots{} (part 4)}
\textbf{Recovery of the standard-Bayes corner.} The KL/NLL client limits
recover Bayes, while the separate zero-robustness server identity
returns the project log-linear pool (eq.~11,
eq.~6, eq.~10). Under the
categorical bridge, that pool specializes only the Eq.
\end{frame}

\begin{frame}{The specific bridge added\ldots{} (part 5)}
7 message-combination term (Friston et al. 2024).
\end{frame}

\begin{frame}{The specific bridge added\ldots{} (part 6)}
\textbf{Additional executed studies.} The source simulations motivate
mechanism analogues, not protocol recovery checks. Added studies cover
declared contamination (sec.~7.5), active
movement (sec.~14.7), two- and three-level POMDPs
(sec.~14.9,
\end{frame}

\begin{frame}{The specific bridge added\ldots{} (part 7)}
sec.~14.11), finite acuity-by-colony sensitivity, and
finite-grid acuity recovery (sec.~14.13,
sec.~15).
\end{frame}

\begin{frame}{The specific bridge added\ldots{} (part 8)}
\textbf{Declared paired analysis.} The configured comparison uses
matched-pairs Wilcoxon tests (Wilcoxon 1945; Fay and Proschan 2010),
Benjamini--Hochberg FDR adjustment (Benjamini and Hochberg 1995),
percentile-bootstrap intervals (Efron and Tibshirani 1993),
\end{frame}

\begin{frame}{The specific bridge added\ldots{} (part 9)}
rank/effect-size caveats (Nakagawa and Cuthill 2007), and an
observed-effect planning approximation. These additions are project
analyses, not claims about the source paper.
\end{frame}

\begin{frame}[fragile]{The specific bridge added\ldots{} (part 10)}
\textbf{Objective-backed server control.}
\texttt{variational\_aggregate} descends its stated free energy
(eq.~12); a converged fixed point is
coordinatewise stationary. The rule has a raw effective-weight bound and
a finite redescending diagnostic (fig.~17).
\end{frame}

\begin{frame}{The specific bridge added\ldots{} (part 11)}
Whether an equally defensible objective has the sharper reverse-KL
heuristic as its minimizer remains open.
\end{frame}

\begin{frame}{The specific bridge added\ldots{} (part 12)}
The contribution is the bridge: an explicit connection between
active-inference belief consensus and robust federated generalized
Bayes,
\end{frame}

\begin{frame}{The specific bridge added\ldots{} (part 13)}
with the standard pool recovered exactly at the corner and the
additional studies (the contamination sweep plus several structural
extensions) showing how the bridge behaves beyond the Friston et
al.~baseline.
\end{frame}

\end{document}
